\documentclass[11pt]{article}
\usepackage{fontspec}
\usepackage{xeCJK}
\setCJKmainfont[ItalicFont=STSongti-SC-Bold]{Songti SC}
\usepackage{xurl}
\usepackage[a4paper,margin=25mm]{geometry}
\usepackage{amsmath,amssymb,amsthm,mathtools}
\usepackage{xcolor}
\usepackage{booktabs,array,enumitem,fancyhdr,hyperref}
\hypersetup{colorlinks=true,linkcolor=blue!50!black,urlcolor=blue!50!black}
\setlength{\emergencystretch}{2em}
\setlength{\parskip}{0.3em}
\linespread{1.16}
\allowdisplaybreaks
\newtheorem{theorem}{定理}[section]
\newtheorem{proposition}[theorem]{命题}
\newtheorem{lemma}[theorem]{引理}
\theoremstyle{definition}
\newtheorem{definition}[theorem]{定义}
\theoremstyle{remark}
\newtheorem{remark}[theorem]{说明}
\newcommand{\dd}{\,\mathrm d}
\newcommand{\ii}{\mathrm i}
\newcommand{\Rea}{\operatorname{Re}}
\newcommand{\Ima}{\operatorname{Im}}
\newcommand{\Li}{\operatorname{Li}}
\newcommand{\supp}{\operatorname{supp}}
\pagestyle{fancy}
\fancyhf{}
\fancyhead[L]{\small 学习笔记}
\fancyhead[R]{\small 初稿}
\fancyfoot[C]{\thepage}
\setlength{\headheight}{15pt}
\author{}
\date{2026年10月22日初稿}
\makeatletter
\renewcommand{\maketitle}{%
  \thispagestyle{plain}%
  \begin{center}
    {\LARGE\bfseries \@title\par}
    \vspace{0.6em}
    {\small \@date\par}
  \end{center}
  \vspace{0.6em}}
\makeatother
\newcommand{\draftnotice}{\noindent{\small
本文为供逐篇审阅的学习笔记初稿，不是新增研究成果。
文中区分直接证明与引用结果；正文不依赖本文。
本轮已作公式、条件及排版初步核查，尚待作者逐句审定。}\par\medskip}

\XeTeXlinebreaklocale "zh"
\renewcommand{\proofname}{证明}
\date{2026年10月24日初稿}
\title{黎曼猜想为何有众多等价形式}
\begin{document}
\maketitle

\noindent{\small 本篇整理 RH 的若干等价形式，并说明它们背后的共同机制。
较易方向给出证明，较难方向注明出处。}\par\medskip

\section{问题}
RH 的等价表述遍及素数分布、算术函数、Farey 序列、泛函分析、热方程与正性条件。
本篇讨论两个问题：在 RH 尚未解决时，“等价”是什么意思；
以及为什么如此不同的命题会与零点位置等价。

\section{“等价”的含义}
在经典逻辑中，若 RH 为真，则 RH 与任何定理都等价；若为假，则与任何假命题都等价。
因此“$X$ 与 RH 等价”作为真值陈述没有信息量，它的实际含义是：
存在 $X\Rightarrow\mathrm{RH}$ 与 $\mathrm{RH}\Rightarrow X$ 的证明，且这些证明都不依赖于确定 RH 的真值。

更有内容的表述是参数之间的恒等式。令
\[
 \Theta=\sup\{\Rea\rho:\ \zeta(\rho)=0,\ 0<\Rea\rho<1\}\in[\tfrac12,1].
\]
RH 即 $\Theta=\tfrac12$。许多等价形式实际上来自形如“某量的增长指数等于 $\Theta$”的无条件定理，
RH 只是其中取 $\Theta=\tfrac12$ 的特例。下一节给出这种机制的两个例子。

\section{机制：Mellin 变换把零点变为增长指数}
\subsection{\texorpdfstring{$\psi$ 的情形}{psi 的情形}}
对 $\Rea s>1$，由 Abel 求和，
\[
 -\frac{\zeta'}{\zeta}(s)=s\int_1^\infty\psi(x)x^{-s-1}\dd x,\qquad
 \frac{s}{s-1}=s\int_1^\infty x\cdot x^{-s-1}\dd x,
\]
故
\[
 -\frac{\zeta'}{\zeta}(s)-\frac{s}{s-1}=s\int_1^\infty\bigl(\psi(x)-x\bigr)x^{-s-1}\dd x .  \tag{3.1}
\]
\begin{proposition}\label{prop:psi}
若 $\psi(x)-x=O(x^\theta)$，$\tfrac12\le\theta<1$，则 $\zeta$ 在 $\Rea s>\theta$ 内无零点。
\end{proposition}
\begin{proof}
在 $\Rea s\ge\theta+\delta$ 上，
$|(\psi(x)-x)x^{-s-1}|\le Cx^{-1-\delta}$。
在任意紧集上，参数 $s$ 有界，求导所产生的固定次 $\log x$ 因子也仍然可积，
故右端在 $\Rea s>\theta$ 内解析。左端在 $s=1$ 处两个极点的留数均为 $1$，相互抵消。
若 $\zeta(\rho)=0$，$\Rea\rho>\theta$，则 $-\zeta'/\zeta$ 在 $\rho$ 处有留数为负重数的极点，
左端不解析，矛盾。
\end{proof}
反方向由显式公式得到：$\psi(x)-x=O(x^\Theta\log^2x)$（Ingham, \emph{The Distribution of Prime Numbers}, 1932）。
两者合起来给出无条件恒等式
\[
 \inf\{\theta:\ \psi(x)-x=O(x^\theta)\}=\Theta .
\]
von Koch（1901）的等价形式 $\pi(x)=\operatorname{li}(x)+O(\sqrt x\log x)$ 是这一恒等式在 $\Theta=\tfrac12$ 时的精细化。

\subsection{Mertens 函数的情形}
同样，对 $\Rea s>1$，$M(x)=\sum_{n\le x}\mu(n)$ 满足
\[
 \frac1{\zeta(s)}=s\int_1^\infty M(x)x^{-s-1}\dd x .                   \tag{3.2}
\]
\begin{proposition}\label{prop:mertens}
若 $M(x)=O(x^\theta)$，$\theta<1$，则 $\zeta$ 在 $\Rea s>\theta$ 内无零点。
\end{proposition}
\begin{proof}
右端在 $\Rea s>\theta$ 内解析，从而 $1/\zeta$ 在该区域解析；$\zeta$ 的零点将是 $1/\zeta$ 的极点。
\end{proof}
反方向（RH $\Rightarrow M(x)=O(x^{1/2+\varepsilon})$）需要在 RH 下估计 $1/\zeta$ 在竖线上的增长，
见 Titchmarsh, \emph{The Theory of the Riemann Zeta-Function}, \S14.25。由此得到 Littlewood（1912）的等价形式。

\subsection{一般原则}
上面两例的结构是：某个算术量的 Mellin 变换（或 Dirichlet 级数）
等于 $\zeta$ 的某个组合（$1/\zeta$、$\zeta'/\zeta$ 等），零点成为它的奇点；
而 Mellin 变换的解析区域由原函数的增长指数决定。
于是“零点不超出 $\Rea s=\tfrac12$”与“误差项不超出 $x^{1/2+\varepsilon}$”相互转化。
若相关零点确实成为不被消去的奇点，并且反方向所需的增长估计和反演步骤也得到证明，
就能建立这类等价形式。仅在生成函数中出现 $\zeta$，还不足以保证两个方向都成立。
这解释了第一类等价形式为何如此之多。

\section{一个泛函分析的例子：Nyman--Beurling 判据}
对 $0<\theta\le1$，令 $\rho_\theta(x)=\{\theta/x\}-\theta\{1/x\}$。
当 $x>1$ 时 $\{\theta/x\}=\theta/x=\theta\{1/x\}$，故 $\rho_\theta$ 支在 $(0,1]$ 上，且有界。
\begin{theorem}[Nyman 1950，Beurling 1955，引用]
RH 成立当且仅当 $\{\rho_\theta:0<\theta\le1\}$ 的线性组合在 $L^2(0,1)$ 中稠密
（等价地，其闭包包含常函数 $\chi=\mathbf 1_{(0,1)}$）。
\end{theorem}
\begin{proposition}\label{prop:nb}
若常函数 $\chi$ 属于上述线性组合的 $L^2(0,1)$ 闭包，则 $\zeta$ 在 $\Rea s>\tfrac12$ 内无零点。
\end{proposition}
\begin{proof}
由第 4 篇的 Abel 表示，在 $0<\Rea s<1$ 内，
\[
 \int_1^\infty\{u\}u^{-s-1}\dd u=\frac1{s-1}-\frac{\zeta(s)}s,
 \qquad
 \int_0^1\{u\}u^{-s-1}\dd u=\frac1{1-s}.
\]
两式相加得 $\int_0^\infty\{u\}u^{-s-1}\dd u=-\zeta(s)/s$。
作代换 $x=\theta/u$ 与 $x=1/u$，得
\[
 \int_0^1\rho_\theta(x)x^{s-1}\dd x=\frac{\zeta(s)}{s}\bigl(\theta-\theta^s\bigr),\qquad 0<\Rea s<1 .
\]
设有限线性组合 $f_n=\sum_jc_{n,j}\rho_{\theta_{n,j}}$ 在 $L^2(0,1)$ 中收敛于 $\chi$。
当 $\Rea s>\tfrac12$ 时 $x^{s-1}\in L^2(0,1)$，由 Cauchy--Schwarz 不等式，
\[
 \int_0^1f_n(x)x^{s-1}\dd x\longrightarrow\int_0^1x^{s-1}\dd x=\frac1s .
\]
若 $\zeta(\rho)=0$，$\tfrac12<\Rea\rho<1$，左端在 $s=\rho$ 处恒为零，而右端为 $1/\rho\ne0$，矛盾。
$\Rea s\ge1$ 处 $\zeta$ 本无零点。
\end{proof}
这里的机制与第 3 节相同：Mellin 变换把伸缩 $x\mapsto x/\theta$ 变为乘以 $\theta^s$，
把 $\{1/x\}$ 变为 $\zeta(s)/s$；零点成为逼近的障碍，而 $L^2$ 的配对恰在 $\Rea s>\tfrac12$ 时有效。
临界线 $\tfrac12$ 在此作为 $L^2$ 的指标出现。

\section{等价形式的分类}
下表按机制列出若干等价形式。“$\varepsilon$”表示对每个 $\varepsilon>0$ 成立。
\begin{center}
\small
\renewcommand{\arraystretch}{1.25}
\begin{tabular}{>{\raggedright\arraybackslash}p{0.2\textwidth}>{\raggedright\arraybackslash}p{0.5\textwidth}>{\raggedright\arraybackslash}p{0.2\textwidth}}
\toprule
机制 & 等价形式 & 出处\\
\midrule
素数分布误差项 & $\pi(x)=\operatorname{li}(x)+O(\sqrt x\log x)$；$\psi(x)=x+O(x^{1/2+\varepsilon})$ & von Koch 1901\\
算术函数和 & $M(x)=O(x^{1/2+\varepsilon})$ & Littlewood 1912\\
 & Redheffer 矩阵 $A_n$ 满足 $\det A_n=M(n)$，故 RH $\iff\det A_n=O(n^{1/2+\varepsilon})$ & Redheffer 1977\\
除数和不等式 & $\sigma(n)<e^\gamma n\log\log n$，$n>5040$ & Robin 1984\\
 & $\sigma(n)\le H_n+e^{H_n}\log H_n$，$n\ge1$ & Lagarias 2002\\
Farey 序列 & $\sum_{\nu}\bigl|f_\nu-\nu/\Phi(N)\bigr|=O(N^{1/2+\varepsilon})$ & Franel, Landau 1924\\
幂级数增长 & $\sum_{k\ge1}\dfrac{(-x)^k}{(k-1)!\,\zeta(2k)}=O(x^{1/4+\varepsilon})$ & Riesz 1916\\
函数逼近 & Nyman--Beurling 判据（第 4 节） & Nyman 1950, Beurling 1955\\
正性 & 显式公式给出的二次型在适当检验函数类上非负 & Weil 1952\\
 & $\lambda_n\ge0$ 对一切 $n$ 成立（学习笔记第 10 篇） & Li 1997\\
导函数零点 & $\zeta'$ 在 $0<\Rea s<\tfrac12$ 内无零点 & Speiser 1934\\
热方程形变 & de Bruijn--Newman 常数 $\Lambda\le0$（由 Rodgers--Tao，即 $\Lambda=0$） & de Bruijn 1950, Newman 1976\\
\bottomrule
\end{tabular}
\end{center}
表中 $f_\nu$ 为 $N$ 阶 Farey 序列在 $(0,1]$ 中的第 $\nu$ 项，$\Phi(N)$ 为其项数；Riesz 判据中级数前的符号约定各文献不一，不影响结论。
Redheffer 矩阵的元素为 $A_{ij}=1$（当 $j=1$ 或 $i\mid j$），其余为 $0$，其中 $1\le i,j\le n$。
Weil 一项仅标出机制，使用时必须按原判据指定检验函数类、卷积与对合，
不能将表中概述直接当作任意二次型的非负定理。

\section{讨论}
\begin{enumerate}
 \item 一大类等价形式经由相关机制产生：零点通过 Mellin 变换或显式公式与增长指数相联系。
 本文已在 $\psi$ 的情形说明增长指数恒等式，在 $M$ 的情形说明一个方向并引用另一个方向。
 Robin 与 Lagarias 的精确不等式还需要额外的算术分析，不能从上述一般机制直接推出。
 \item 第二种机制是正性：Weil 与 Li 的判据把零点位置转化为一族二次型或一列实数的非负性。
 第三种是逼近或谱的语言（Nyman--Beurling），其中临界线以 $L^2$ 的指标出现。
 \item 等价改写一般不降低问题的难度，但提供不同的工具与不同的逻辑形式：
 Lagarias 判据的严格反例若存在，可以由有限计算认证（学习笔记第 13 篇）；
 一般可计算实数间的非严格比较不自动给出保证终止的判定过程。
 Li 判据适合数值观察；de Bruijn--Newman 常数允许讨论“离 RH 有多远”。
 \item 等价形式之多也反映了 $\zeta$ 在数论中的位置：许多算术部分和通过其生成函数与零点相联系。
 但不同问题还可能涉及其他 $L$ 函数、额外奇点或新的估计条件，
 不能断言所有乘性函数的误差都由 $\zeta$ 零点单独决定。
\end{enumerate}

\section*{参考}
\begin{itemize}
 \item E.\ C.\ Titchmarsh, \emph{The Theory of the Riemann Zeta-Function}, 2nd ed., 1986, Ch.~II, XIV.
 \item A.\ E.\ Ingham, \emph{The Distribution of Prime Numbers}, Cambridge, 1932.
 \item P.\ Borwein, S.\ Choi, B.\ Rooney, A.\ Weirathmueller (eds.), \emph{The Riemann Hypothesis: A Resource for the Afficionado and Virtuoso Alike}, Springer, 2008（汇集了多数等价形式的原始文献）。
 \item G.\ Robin, Grandes valeurs de la fonction somme des diviseurs et hypothèse de Riemann, \emph{J.\ Math.\ Pures Appl.} 63 (1984).
 \item J.\ C.\ Lagarias, An elementary problem equivalent to the Riemann hypothesis, \emph{Amer.\ Math.\ Monthly} 109 (2002).
 \item A.\ Beurling, A closure problem related to the Riemann zeta-function, \emph{PNAS} 41 (1955).
 \item X.-J.\ Li, The positivity of a sequence of numbers and the Riemann hypothesis, \emph{J.\ Number Theory} 65 (1997).
\end{itemize}
\end{document}
